\pagebreak

\section{Floating Point Library - float.3pl}\label{chap:fp}

    \index{library!float.3pl}
    This library provides floating point modules, procedures and functions for
    all Xilinx FPGAs.
    
    This library must be incorporated by using a {\bf module} statement, i.e.
    \begin{lstlisting}[gobble=4, language=threepl]
       module "float.3pl"
    \end{lstlisting}
    
    This library was developed for a specific project, but it was soon determined
    that fixed point operations were far better suited and used far fewer
    logic resources. Hence a fixed point target type was incorporated fully
    into 3PL and this floating point library has never been used in earnest.
    It passed early testing but probably needs much more development.



    \subsection{library modules}


        \subsubsection{fadd - floating add}\label{sec:lmfadd}

            {\bf Library: float.3pl}

            {\bf Prototype:} \\
            \vsp
            \begin{lstlisting}[gobble=12, language=threepl]
               global module
               fadd (queue(out) <- value(in1), value(in2)) {}
            \end{lstlisting}

            {\bf Output parameters:} \\
            {\bf out} is the output queue variable.

            {\bf Input parameters:} \\
            {\bf in1} is the first input value. \\
            {\bf in1} is the second input value.

            {\bf description:} \\
            This module adds the two input floating point values and assigns
            the sum to the output queue. The input values and the output queue
            must be type {\bf float}. The input exponent and mantissa widths
            may differ from each other and from those of the output format. The
            module will assign data to the output queue when both inputs and the
            output are available.
            
            The latency through this module is $5 + \lceil \log_2 m \rceil + r$
            where $m$ is the number of bits in the mantissa and $r$ is 1 if
            either of the input formats differs to the output format and 0
            otherwise.
            
            Overflow and underflow are not handled.


        \subsubsection{fdiv - floating divide}\label{sec:lmfdiv}

            {\bf Library: float.3pl}

            {\bf Prototype:} \\
            \vsp
            \begin{lstlisting}[gobble=12, language=threepl]
               global module
               fdiv (queue(out) <- value(in1), value(in2)) {}
            \end{lstlisting}

            {\bf Output parameters:} \\
            {\bf out} is the output queue variable.

            {\bf Input parameters:} \\
            {\bf in1} is the first input value. \\
            {\bf in1} is the second input value.

            {\bf description:} \\
            This module divides the first input floating point value by the
            second and assigns the quotient to the output queue. The input
            values and the output queue must be type {\bf float}. The input
            exponent and mantissa widths may differ from each other and from
            those of the output format. The module will assign data to the
            output queue when both inputs and the output are available.        
            
            Overflow and underflow are not handled.
            
            {\bf {\large This module has not yet been written}}


        \subsubsection{fmul - floating multiply}\label{sec:lmfmul}

            {\bf Library: float.3pl}

            {\bf Prototype:} \\
            \vsp
            \begin{lstlisting}[gobble=12, language=threepl]
               global module
               fmul (queue(out) <- value(in1), value(in2)) {}
            \end{lstlisting}

            {\bf Output parameters:} \\
            {\bf out} is the output queue variable.

            {\bf Input parameters:} \\
            {\bf in1} is the first input value. \\
            {\bf in1} is the second input value.

            {\bf description:} \\
            This module multiplies the two input floating point values and
            assigns the product to the output queue. The input values and the
            output queue must be type {\bf float}. The input exponent and
            mantissa widths may differ from each other and from those of the
            output format. The module will assign data to the output queue when
            both inputs and the output are available.
            
            The latency through this module is $2 + \lceil \log_2 m \rceil + r$
            where $m$ is the number of bits in the mantissa and $r$ is 1 if
            either of the input formats differs to the output format and 0
            otherwise.
            
            Overflow and underflow are not handled.


        \subsubsection{fneg - negate a floating point value}\label{sec:lmfneg}

            {\bf Library: float.3pl}

            {\bf Prototype:} \\
            \vsp
            \begin{lstlisting}[gobble=12, language=threepl]
               global module
               fneg (queue(out) <- value(in)) {}
            \end{lstlisting}

            {\bf Output parameters:} \\
            {\bf out} is the output queue variable.

            {\bf Input parameters:} \\
            {\bf in} is the input value.

            {\bf description:} \\
            This module negates the input floating point value and assigns the
            result to the output queue. The input value and the output queue must
            be type {\bf float}. The input exponent and mantissa widths may
            differ from those of the output format. The module will assign data
            to the output queue when both input and output are available.

            The latency through this module is 2 if the input format
            differs to that of the output and 1 otherwise.
            
            This module simply changes the sign bit, including the


        \subsubsection{fsub - floating subtract}\label{sec:lmfsub}

            {\bf Library: float.3pl}

            {\bf Prototype:} \\
            \vsp
            \begin{lstlisting}[gobble=12, language=threepl]
               global module
               fsub (queue(out) <- value(in1), value(in2)) {}
            \end{lstlisting}

            {\bf Output parameters:} \\
            {\bf out} is the output queue variable.

            {\bf Input parameters:} \\
            {\bf in1} is the first input value. \\
            {\bf in1} is the second input value.

            {\bf description:} \\
            This module subtracts the second input floating point value
            from the first and assigns the difference to the output queue.
            The input values and the output queue must be type {\bf float}.
            The input exponent and mantissa widths may differ from each other
            and from those of the output format.
            The module will assign data to the output queue when both
            inputs and the output are available.
            
            The latency through this module is $5 + \lceil \log_2 m \rceil + r$
            where $m$ is the number of bits in the mantissa and $r$ is 1 if
            either of the input formats differs to the output format and 0
            otherwise.
            
            Overflow and underflow are not handled.


        \subsubsection{itof - convert an integer to a floating point value}\label{sec:lmitof}

            {\bf Library: float.3pl}

            {\bf Prototype:} \\
            \vsp
            \begin{lstlisting}[gobble=12, language=threepl]
               global module
               itof (queue(out) <- value(in)) {}
            \end{lstlisting}

            {\bf Output parameters:} \\
            {\bf out} is the output queue variable.

            {\bf Input parameters:} \\
            {\bf in} is the input queue variable.

            {\bf description:} \\
            This module converts the integer in the input argument to a floating
            point value which is assigned to the output queue.
            The input value must be type {\bf int} or {\bf uint} and the
            output queue must be type {\bf float}.
            The module will assign data to the output queue when both
            input and output arguments are available.
            
            The latency through this module is $1 + \lceil \log_2 m \rceil$
            where $m$ is the number of bits in the mantissa.
            
            Overflow is not handled, but is very unlikely to occur in converting
            integer to floating point.





    \subsection{library procedures}
    
        None.




    \subsection{library functions}
    

        None.


    \subsection{library classes}
    
        None.
